\section{Analysis of Primary Studies and Architectural Gaps}
\label{sec:analysis-studies}

The synthesis of the $N = 10$ primary studies reveals critical insights into the current state of the art in cloud-native observability. Table~\ref{tab:rq_mapping} maps the included primary studies against the four review research questions (\textbf{RQ01}--\textbf{RQ04}), detailing their publication year, venue type, telemetry scope, and quality assessment (QA) score, while Table~\ref{tab:related_works} synthesizes the specific architectural limitations and future research directions extracted from each work.

\begin{table*}[!t]
\centering
\caption{Primary Studies Included in the Rapid Review, Quality Assessment, and RQ Mapping ($N = 10$)}
\label{tab:rq_mapping}
\footnotesize
\renewcommand{\arraystretch}{1.15}
\begin{tabularx}{\textwidth}{@{} p{2.3cm} l X X c c c c c @{}}
\toprule
\textbf{Primary Study} & \textbf{Year} & \textbf{Venue / Publication Type} & \textbf{Telemetry / Primary Scope} & \textbf{QA} & \textbf{RQ1} & \textbf{RQ2} & \textbf{RQ3} & \textbf{RQ4} \\
\midrule
Kosińska et al. \cite{kosinska_2023} & 2023 & IEEE Access (Journal) & Metrics, Logs, Traces (Full-stack state-of-the-art) & 4.0 & \checkmark & \checkmark & \checkmark & \checkmark \\
\addlinespace
Usman et al. \cite{usman_2022} & 2022 & IEEE Access (Journal) & Metrics, Logs, Traces (Edge \& Container microservices) & 4.0 & \checkmark & \checkmark & \checkmark & \checkmark \\
\addlinespace
Faseeha et al. \cite{faseeha_2025} & 2025 & IEEE Access (Journal) & Metrics, Logs, Traces, Events (Deployment paradigms) & 4.0 & \checkmark & \checkmark & \checkmark & \checkmark \\
\addlinespace
Gomes et al. \cite{gomes_2025} & 2025 & Computing / Springer (Journal) & Metrics, Logs, Traces (Auto-scaling \& Taxonomies) & 4.0 & \checkmark & & \checkmark & \checkmark \\
\addlinespace
Han et al. \cite{han_2024} & 2024 & IEEE Trans. Serv. Comput. (Journal) & Metrics, Logs, Traces (Root Cause Analysis \& GAT) & 4.0 & \checkmark & \checkmark & & \checkmark \\
\addlinespace
Deng et al. \cite{deng_2024} & 2024 & Proc. of the IEEE (Journal) & Metrics, Traces, Events (Service Lifecycle \& DevEx) & 3.5 & & \checkmark & & \checkmark \\
\addlinespace
Araujo et al. \cite{araujo_2024} & 2024 & arXiv / ACM DL (Preprint/Workshop) & Metrics, Logs, Traces (Fog Computing Telemetry Lifecycle) & 3.5 & & \checkmark & & \checkmark \\
\addlinespace
Hartikainen \cite{hartikainen_2024} & 2024 & Häme UAS (Master's Thesis) & Metrics, Logs, Traces, Events (Kubernetes LGTM / OTel) & 3.5 & \checkmark & & & \\
\addlinespace
Bruno et al. \cite{zanatta_23} & 2023 & JISA / SBC (Journal) & Metrics, Logs (Beyond-5G CNOT Benchmark \& Overhead) & 3.0 & & & & \checkmark \\
\addlinespace
Wei et al. \cite{wei_2025} & 2025 & IEEE/ACM ICSA (Conference) & Metrics, Traces (Quality-to-Tool Mapping) & 3.0 & & & & \checkmark \\ 
\bottomrule
\end{tabularx}
\end{table*}

\subsection{Essential Architectural Components (RQ01)}
The literature establishes that an effective cloud-native observability architecture requires coordinated capabilities across multiple functional layers \cite{kosinska_2023, usman_2022, faseeha_2025}. The core data foundation comprises the telemetry tuple $O = \{M, L, T, E\}$ (metrics, logs, distributed traces, and platform events). However, primary studies converge on the fundamental principle that collecting these pillars in isolated silos yields limited diagnostic value; high-value observability demands cross-pillar semantic correlation at the points of ingestion, storage, and query analysis \cite{gomes_2025, araujo_2024}.

Functionally, the architecture must support the end-to-end Observability Data Life Cycle (ODLC) \cite{araujo_2024, han_2024}: (1)~\textit{Collection and Instrumentation}, acquiring raw telemetry via language SDKs, sidecars, daemon exporters, and kernel probes; (2)~\textit{Processing and Pipeline Buffering}, performing filtering, label enrichment, batching, and dynamic sampling; (3)~\textit{Specialized Storage}, distributing telemetry into specialized time-series, chunked-index, or columnar object databases; and (4)~\textit{Correlation and Diagnosis}, synthesizing multidimensional queries, service dependency graphs, and root-cause localization workflows.

\subsection{Reference Architectures and Frameworks (RQ02)}
A major finding of the review is the noticeable absence of formal, standardized reference architectures in the scientific literature specifically engineered and empirically validated for Kubernetes \cite{araujo_2024, kosinska_2023}. Although academic works propose high-level conceptual lifecycle models \cite{araujo_2024, han_2024}, production environments rely largely on community-driven \textit{de facto} deployment patterns structured across four tiers:
\begin{enumerate}
    \item \textbf{Agent Tier:} Collectors deployed as cluster-wide Kubernetes DaemonSets (e.g., OpenTelemetry Collector, Fluent Bit, Node Exporter), pod sidecars, or kernel probes (eBPF) to intercept runtime signals.
    \item \textbf{Aggregation and Routing Tier:} Centralized collector pipelines that ingest heterogeneous telemetry protocols, normalize metadata schemas, apply batching, and route payloads to designated storage backends.
    \item \textbf{Specialized Datastore Tier:} Decoupled storage engines tailored to distinct telemetry shapes---such as time-series databases (Prometheus) for numeric metrics, label-indexed chunk storage (Grafana Loki) for logs, and object-backed distributed trace stores (Jaeger, Grafana Tempo).
    \item \textbf{Query and Correlation Tier:} Unified visualization and query interfaces (primarily Grafana) providing single-pane-of-glass dashboards, exemplar-based drill-down, and programmatic HTTP query endpoints for automated Root Cause Analysis.
\end{enumerate}

\subsection{Open-Source Toolchains and Interoperability (RQ03)}
The review identifies a mature, yet operationally fragmented, open-source toolchain across the CNCF ecosystem. OpenTelemetry (OTel) has consolidated itself as the leading vendor-neutral framework for telemetry generation and distributed trace propagation via the OpenTelemetry Protocol (OTLP) \cite{kosinska_2023, faseeha_2025}. Nevertheless, OTel is not yet a frictionless universal standard across all storage dimensions.

For metric scraping and alerting, Prometheus and its OpenMetrics exposition format remain the dominant standard in Kubernetes deployments \cite{usman_2022, faseeha_2025}. Bidirectional interoperability between OpenTelemetry and Prometheus still presents non-trivial engineering challenges, including lossy metric conversions, differences in histogram bucket representations, and label naming conventions \cite{hartikainen_2024}. In the logging domain, Grafana Loki and Fluentd/Fluent Bit lead Kubernetes log management, whereas Jaeger and Grafana Tempo dominate distributed tracing. Despite the maturity of individual tools, primary studies emphasize that these components are rarely plug-and-play: engineering teams face high cognitive load and operational overhead in configuring collectors, enforcing label parity across pillars, and preventing resource contention on worker nodes \cite{kosinska_2023, usman_2022}.

\subsection{Synthesis of Research Gaps (RQ04)}

\begin{table*}[!ht]
\centering
\caption{Analysis of Primary Studies: Research Gaps, Limitations, and Future Works ($N = 10$)}
\label{tab:related_works}
\footnotesize
\begin{tabular}{p{0.18\textwidth} p{0.38\textwidth} p{0.38\textwidth}}
\toprule
\textbf{Primary Study} & \textbf{Research Gaps and Limitations Identified} & \textbf{Future Directions and Proposed Work} \\ \midrule

Usman et al. (2022) \cite{usman_2022} &
Absence of formal reference architectures; severe operational complexity arising from stack fragmentation; telemetry scalability and multimodal correlation bottlenecks. & 
Intelligent automation via AI/ML (AIOps); integrated runtime security (SecOps); enhancement of developer experience (DevEx); unified observability platforms. \\ \addlinespace

Araujo et al. (2024) \cite{araujo_24} &
Lack of holistic reference architectures; severe CPU/memory overhead constraints in resource-constrained edge/fog nodes; high device heterogeneity. & 
Lightweight AI/ML for decentralized anomaly detection (TinyML); standardized interoperability protocols; adaptive low-overhead observability frameworks. \\ \addlinespace

Bruno et al. (2023) \cite{zanatta_23} &
Biased focus on metrics over logs and distributed traces in B5G research; significant memory and CPU overhead of standard CNOT stacks; lack of standardized log schemas. & 
Integration of end-to-end distributed tracing; scalability analysis of COTS ML detectors; non-intrusive telemetry for disaggregated Open RAN components. \\ \addlinespace

Kosińska et al. (2023) \cite{kosinska_2023} &
Lack of unified terminology across layers; insufficient practical automation via AI/ML (AIOps gap); poor adaptation of tools for edge and resource-limited topologies. & 
Development of formal observability maturity models; deeper eBPF integration for non-intrusive visibility; alignment with FinOps and GreenOps principles. \\ \addlinespace

Faseeha et al. (2025) \cite{faseeha_2025} &
Conceptual confusion between traditional monitoring and modern observability pillars; lack of practical multi-pillar approaches; difficulties tracing requests in ephemeral containers. & 
Evolution of frameworks for modern deployment paradigms (service mesh, event-driven); automated telemetry analytics; systematic minimization of runtime overhead. \\ \addlinespace

Han et al. (2024) \cite{han_2024} &
Inconsistent problem formulations for root cause analysis (RCA); high cloud-native heterogeneity across traces, metrics, and logs; high correlation compute costs. & 
Scalable multimodal data correlation algorithms; graph attention networks (GAT) for dynamic topologies; automated mitigation workflows. \\ \addlinespace

Gomes et al. (2025) \cite{gomes_2025} &
Difficulty in correlating multimodal telemetry data; inadequacy of traditional monitoring for microservices; complex root cause localization for autoscaling. & 
Advanced ML-based resource management; automated diagnostic mechanisms coupled with autoscaling controllers; low-overhead real-time collection. \\ \addlinespace

Wei et al. (2025) \cite{wei_2025} &
Lack of a systematic bridge between desirable system qualities and concrete cloud-native tools; limited initial empirical validation across production settings. &
Interfaces to continuously evolve quality-to-tool mappings; comprehensive industrial validation beyond initial expert consensus. \\ \addlinespace

Deng et al. (2024) \cite{deng_2024} &
Missing unified research roadmap for cloud-native lifecycles; heavy cognitive load for developers (DevEx); lack of real-time security observability. & 
Integration of SecOps with performance telemetry; improvement of Developer Experience (DevEx); observability optimization for serverless and green computing. \\ \addlinespace

Hartikainen (2024) \cite{hartikainen_2024} &
High integration complexity in multi-tenant Kubernetes clusters; configuration divergence between OTel and native collectors; high storage ingestion costs. &
Production benchmark evaluations of LGTM vs OTel collector architectures; automated retention and adaptive sampling pipelines for Kubernetes platforms. \\ \bottomrule
\end{tabular}
\end{table*}

As synthesized in Table~\ref{tab:related_works}, the state of the art exposes three pervasive research gaps that directly motivate this thesis:
\begin{enumerate}
    \item \textbf{Stack Fragmentation and Absence of Semantic Correlation Schemas:} Disconnected storage silos prevent automated end-to-end root-cause localization across disaggregated microservices \cite{usman_2022, gomes_2025};
    \item \textbf{Unbounded Computational and Memory Overhead:} Full-fidelity telemetry collection and heavy full-text indexing (e.g., Elasticsearch/Cassandra stacks) impose severe CPU, memory, and I/O overhead in resource-constrained Kubernetes nodes \cite{zanatta_23, araujo_2024}; and
    \item \textbf{Telemetry Overload and Diagnostic Latency during RCA:} The sheer volume of unindexed logs and cascading trace spans overwhelms operators during complex incident response, inflating the Mean Time to Diagnose ($MTTDiag$) \cite{han_2024, deng_2024}.
\end{enumerate}
